% Included from main.tex after \appendix.
% This file intentionally does not redefine the document class or packages.



\subsection{A Note on the Energy Eigenvalue in the Corrected Schrödinger Equation}
\label{app:corrected_schrodinger_eigenvalue}

The relativistically corrected Schrödinger equation is commonly written as



In particular, the physical time dependence of a stationary relativistic
state remains
\begin{equation}
\Psi(\mathbf r,t)
=
\psi(\mathbf r)
\exp\left(-\frac{iE_{\mathrm{tot}}t}{\hbar}\right).
\end{equation}
Thus, $E_{\mathrm{tot}}$ governs the temporal phase, whereas
$p_0^2/(2\gamma m_0)$ is the eigenvalue of the equivalent stationary
Schrödinger-like equation.



\section{Relativistic Hamiltonian and Scalar Wave Equations}
\label{app:relativistic_hamiltonian}

This appendix gives the derivation underlying the two scalar refractive-index
models used in Section~\ref{sec:schrodinger_helmholtz}. A complete relativistic
description begins with the Dirac equation. When spin and magnetic interactions
are not explicitly required, elastic electron diffraction can instead be
described by the positive-energy, spin-independent relativistic Hamiltonian
\cite{fujiwara1961relativistic,kirkland2010advanced}.

\subsection{Minimal Coupling and the Stationary Equation}

The time-dependent wave equation is
\begin{equation}
i\hbar\frac{\partial\Psi(\mathbf r,t)}{\partial t}
=
\hat H\Psi(\mathbf r,t),
\label{eq:app_time_dependent_wave_equation}
\end{equation}
with the minimal-coupling Hamiltonian
\begin{equation}
\hat H
=
\sqrt{
c^2\left(\hat{\mathbf p}-q_e\mathbf A/c\right)^2+m_0^2c^4
}
+
q_eV(\mathbf r),
\label{eq:app_minimal_coupling_hamiltonian}
\end{equation}
where $\hat{\mathbf p}=-i\hbar\nabla$, $\mathbf A$ is the magnetic vector
potential, and $V(\mathbf r)$ is the specimen electrostatic potential in
volts. We use $e>0$ and $q_e=-e$, so the electron potential energy is
\begin{equation}
E_{\mathrm{pot}}(\mathbf r)
=
q_eV(\mathbf r)
=
-eV(\mathbf r).
\label{eq:app_electron_potential_energy}
\end{equation}


This notation distinguishes the electrostatic potential $V$ from the energy
$E_{\mathrm{pot}}$.

Neglecting magnetic interactions and introducing the rest energy $ E_0 = m_0 c^2$ gives
\begin{equation}
\hat H
=
\sqrt{c^2\hat{\mathbf p}^{\,2}+m_0^2c^4}
+
E_{\mathrm{pot}}(\mathbf r) = \sqrt{c^2\hat{\mathbf p}^{\,2}+E_0^2}
+
E_{\mathrm{pot}}(\mathbf r)
\label{eq:app_relativistic_hamiltonian}
\end{equation}

We then look for stationnary states $\Psi(\mathbf r,t)$ that separate the temporal and spatial contributions : 
\begin{equation}
\Psi(\mathbf r,t)
=
\psi(\mathbf r)
\exp\left(-\frac{iE_{\mathrm{tot}}t}{\hbar}\right),
\label{eq:app_stationary_ansatz}
\end{equation}

 The stationary equation is therefore
\begin{equation}
\hat H\psi=E_{\mathrm{tot}}\psi.
\label{eq:app_stationary_equation}
\end{equation}

$E_{\mathrm{tot}}$ is the conserved total relativistic energy, including
the rest energy. As we look for elastic interaction this energy is conserved. In particular, 
in vacuum,
\begin{equation}
E_{\mathrm{tot}}
=
\gamma m_0c^2
=
\sqrt{c^2p_0^2+E_0^2},
\qquad
p_0=hk_0=\frac{h}{\lambda},
\label{eq:app_vacuum_energy_momentum}
\end{equation}. 
If the beam has accelerated through a voltage magnitude $U_{\mathrm{acc}}>0$, the total energy is the sum of the rest and kinetic energies : 

and $E_{\mathrm{kin}}=eU_{\mathrm{acc}}=E_{\mathrm{tot}}-E_0 = (\gamma - 1) E_0$. 

Equation~\eqref{eq:app_vacuum_energy_momentum} can be used to compute the wave length : 
$c^2 p_0^2 = h^2 c^2/\lambda^2 = E_{\mathrm{tot}}^2 - E_0^2 = (E_{\mathrm{kin}}+E_0)^2 - E_0^2$ 


\subsection{Klein--Gordon and Corrected Schr\"odinger Forms}

Rearranging Eq.~\eqref{eq:app_stationary_equation} and squaring gives
\begin{equation}
\left[c^2\hat{\mathbf p}^{\,2}+E_0^2\right]\psi
=
\left[E_{\mathrm{tot}}-E_{\mathrm{pot}}(\mathbf r)\right]^2\psi .
\label{eq:app_kg_before_linearisation}
\end{equation}
A local relattivistic momentum and a local wavelength can be defined : 
\begin{equation}
c^2\hat{\mathbf p}^{\,2}
=
\left[E_{\mathrm{tot}}-E_{\mathrm{pot}}(\mathbf r)\right]^2 - E_0^2 = h^2 c^2 \mathbf k^2
\label{eq:app_kg_before_linearisation2}
\end{equation}

Using $\hat{\mathbf p}^{\,2}=-\hbar^2\nabla^2$ and
$p_0=h k_0$, equation ~\eqref{eq:app_vacuum_energy_momentum} becomes
\begin{equation}
\nabla^2\psi(\mathbf r)
+
(2\pi k_0)^2n_{\mathrm{KG}}^2(\mathbf r)\psi(\mathbf r)
=0,
\label{eq:app_kg_helmholtz_form}
\end{equation}
where
\begin{equation}
n_{\mathrm{KG}}^2(\mathbf r)
=
\frac{
\left[E_{\mathrm{tot}}-E_{\mathrm{pot}}(\mathbf r)\right]^2-E_0^2
} {
E_{\mathrm{tot}}^2-E_0^2
} = \frac{p(\mathbf r)}{p_0}=\frac{k(\mathbf r)}{k_0}
.
\label{eq:app_kg_refractive_index}
\end{equation}
So, $n_{\mathrm{KG}}$ is the ratio of the local relativistic 
momentum $p(r)$  to the vacuum momentum $p_0$ . This connects the Klein--Gordon form to the
electron-optical refractive-index concepts of Glaser and Ehrenberg--Siday
\cite{Glaser1933,Ehrenberg-Siday1949}.

For $|E_{\mathrm{pot}}|\ll E_{\mathrm{tot}}$, as it is the case for electron-beam interaction, 
\begin{equation}
\left[E_{\mathrm{tot}}-E_{\mathrm{pot}}(\mathbf r)\right]^2
\simeq
E_{\mathrm{tot}}^2
-
2E_{\mathrm{tot}}E_{\mathrm{pot}}(\mathbf r).
\label{eq:app_linearised_energy}
\end{equation}
Since $E_{\mathrm{tot}}/c^2=\gamma m_0$, Eq.~\eqref{eq:app_kg_before_linearisation}
then reduces to
\begin{equation}
\nabla^2\psi(\mathbf r)
+
\left[
(2\pi k_0)^2
-
\frac{2\gamma m_0}{\hbar^2}E_{\mathrm{pot}}(\mathbf r)
\right]\psi(\mathbf r)
=0.
\label{eq:app_corrected_schrodinger}
\end{equation} or equivalently , 

\begin{equation}
\left[
-\frac{\hbar^2}{2\gamma m_0}\nabla^2
+
E_{\mathrm{pot}}(\mathbf r)
\right]\psi(\mathbf r)
=
\frac{p_0^2}{2\gamma m_0}\psi(\mathbf r).
\label{eq:app_schrodinger_energy_form}
\end{equation}

Equation~\eqref{eq:app_schrodinger_energy_form} is
therefore a stationary Schrödinger-like reformulation of the relativistic
wave equation, rather than the original relativistic energy-eigenvalue
equation.
This is the relativistically corrected
Schr\"odinger form as deduced by Fujiwara. 
The right-hand side resembles the nonrelativistic kinetic energy
$p_0^2/(2m)$, with the rest mass replaced by the relativistic mass
$\gamma m_0$. It should not, however, be identified with either the total
relativistic energy $E_{\mathrm{tot}}$ or, in general, the kinetic energy
$E_{\mathrm{kin}}=E_{\mathrm{tot}}-E_0$.

Indeed, using
\begin{equation}
\gamma m_0=\frac{E_{\mathrm{tot}}}{c^2},
\qquad
p_0^2c^2=E_{\mathrm{tot}}^2-E_0^2,
\end{equation}
the eigenvalue in Eq.~\eqref{eq:app_schrodinger_energy_form} becomes
\begin{equation}
\frac{p_0^2}{2\gamma m_0}
=
\frac{E_{\mathrm{tot}}^2-E_0^2}{2E_{\mathrm{tot}}}
=
E_{\mathrm{kin}}\frac{E_{\mathrm{tot}}+E_0}{2E_{\mathrm{tot}}}.
\label{eq:app_effective_schrodinger_eigenvalue}
\end{equation}

It approaches $E_{\mathrm{kin}}$ only in the nonrelativistic limit
$E_{\mathrm{kin}}\ll E_0$, for which $E_{\mathrm{tot}}\simeq E_0$.

Equation  xxxx-mettrelabonneREF indicates that the energy appearing in the temporal dependence of the wave function is not $E_{0\mathrm{kin}}$ but $E_{\mathrm{tot}}$. 
The factor $1/2$ arises when the scalar relativistic wave equation is
linearised in the specimen potential and divided by
$2E_{\mathrm{tot}}$. 

\subsection{Interaction Constant and Paraxial Limit}

The conventional TEM interaction constant is
\begin{equation}
\sigma
=
\frac{2\pi\gamma m_0e}{h^2k_0}
=
\frac{2\pi\gamma m_0e\lambda}{h^2}= 2 \pi e  k_0 \frac{\gamma m_0 c^2}{h^2 k_0^2 c^2}=
\frac{2\pi e E_{\mathrm{tot}}}
{hc\sqrt{E_{\mathrm{tot}}^2-E_0^2}}.
\label{eq:app_sigma_definition}
\end{equation}

The potential term in Eq.~\eqref{eq:app_corrected_schrodinger} can therefore be
written as
\begin{equation}
\frac{2\gamma m_0e}{\hbar^2}V(\mathbf r)
=
4\pi k_0\sigma V(\mathbf r),
\label{eq:app_potential_term_sigma}
\end{equation}
giving
\begin{equation}
\nabla^2\psi(\mathbf r)
+
\left[
(2\pi k_0)^2
+
4\pi k_0\sigma V(\mathbf r)
\right]\psi(\mathbf r)
=0.
\label{eq:app_corrected_schrodinger_sigma}
\end{equation}
Factoring out the carrier wave,
\begin{equation}
\psi(\mathbf r)
=
\exp(2\pi i k_0z)\tilde\psi(\mathbf r),
\label{eq:app_carrier_wave}
\end{equation}
and neglecting $\partial^2\tilde\psi/\partial z^2$ gives
\begin{equation}
\frac{\partial\tilde\psi}{\partial z}
=
\frac{i}{4\pi k_0}\nabla_\perp^2\tilde\psi
+
i\sigma V(\mathbf r)\tilde\psi.
\label{eq:app_paraxial_sigma}
\end{equation}
For a slice with projected electrostatic potential
\begin{equation}
V_j(x,y)
=
\int_{z_j}^{z_j+\Delta z}V(x,y,z)\,dz,
\label{eq:app_projected_potential}
\end{equation}
the phase-grating transmission function is
\begin{equation}
t_j(x,y)
=
\exp\left[i\sigma V_j(x,y)\right],
\label{eq:app_transmission_function}
\end{equation}
which is the expression used in the main text.
